\begin{center}\begin{tikzpicture}[font=\small,>=stealth,line join=round]
\draw[gray!40,rounded corners=3pt] (0,0) rectangle (14.2,9.6);
\node at (7.1,9.0) {展开平方的等号，为什么对应两段同向？};

\begin{scope}[shift={(1.9,1.65)},x=1.0cm,y=1.0cm]
\draw[->,gray] (0,0)--(5.0,0) node[right]{水平};\draw[->,gray] (0,0)--(0,5.8) node[above]{竖直};
\draw[gray,dashed] (0,0)--(2,5);
\draw[orange!85!black,very thick] (0,0)--(2,2);\draw[teal!80!black,very thick] (2,2)--(2,5);
\fill (0,0) circle(2pt) (2,2) circle(2pt) (2,5) circle(2pt);
\node[below left] at (0,0){$O$};\node[right=6pt] at (2,2){$C(u,a)$};\node[above right=6pt] at (2,5){$D(u+v,a+b)$};
\node[below=8pt] at (1,1){$U$};\node[right=6pt] at (2,3.5){$V$};
\end{scope}
\node[align=left,anchor=west] at (8.4,4.9){图取 $u=2,v=0,a=2,b=3$\\[12pt]$OC=U,\ CD=V$\\[12pt]$(UV)^2-(uv+ab)^2$\\[8pt]$\qquad=(bu-av)^2\ge0$\\[14pt]$bu=av$ 时两段同向\\[8pt]折点消失，等号成立};

\node[text width=12.8cm,align=center] at (7.1,.65) {这里只用图解释等号的形状；正文的平方展开对正、负的 $u,v$ 都成立。};
\end{tikzpicture}\end{center}
